%*******************************************************************************
%****************************** Second Chapter *********************************
%*******************************************************************************

\chapter{Review of Related Literature}

\ifpdf
    \graphicspath{{Chapter2/Figs/Raster/}{Chapter2/Figs/PDF/}{Chapter2/Figs/}}
\else
    \graphicspath{{Chapter2/Figs/Vector/}{Chapter2/Figs/}}
\fi

Effective patient information management systems (PIMS) are indispensable in modern healthcare settings, improving efficiency, accuracy, and patient care \cite{who2010}. However, university clinics often face unique challenges due to their diverse patient populations, resource constraints, and specific operational needs \cite{nagel2022}. Modernizing the Patient Information System (PIS) at the USC Clinic is crucial to address these challenges and leverage the latest technological advancements to enhance healthcare delivery.

\section{Challenges in University Clinic PIMS}

Research highlights several key challenges in university clinic PIMS that must be addressed during modernization:

\subsection{Incomplete Patient Data}
\begin{enumerate}
    \item Studies have shown that incomplete or inaccurate patient data can lead to medical errors, missed diagnoses, and delayed treatment \cite{chowdhury2022}.
    \item In university settings, capturing comprehensive health information, especially for incoming students, can be particularly challenging due to various factors such as time constraints, limited resources, and varying levels of student engagement \cite{donnelly2022}.
    \item Effective strategies for addressing this issue include integrating student health records with enrollment systems and implementing proactive outreach programs to encourage students to provide their medical information \cite{donnelly2022, wurster2024}.
\end{enumerate}

\subsection{Inefficient Consultation Management}
\begin{enumerate}
    \item Long wait times and inefficient scheduling are common issues in university clinics, leading to patient dissatisfaction and potential delays in care \cite{biya2022}.
    \item Implementing online self-scheduling tools and utilizing telemedicine for minor ailments can help improve patient flow and reduce wait times \cite{byun2022}.
\end{enumerate}

\subsection{Medication Inventory Control}
\begin{enumerate}
    \item Manual inventory management in clinics can be error-prone and lead to stockouts or expired medications \cite{deressa2022}.
    \item Electronic inventory systems, especially those using the First-In, First-Out (FIFO) principle, have been shown to optimize medication usage, reduce waste, and improve cost-effectiveness \cite{sohrabi2021}.
\end{enumerate}

\subsection{Daily Treatment Records (DTR) and Reporting}
Paper-based or spreadsheet-based DTRs are inefficient and hinder data analysis. Electronic DTRs integrated within a PIMS can automate reporting, improve data accuracy, and provide valuable insights into clinic utilization and patient outcomes \cite{baumann2018}.

\subsection{Medical Certificate Issuance}
Manual issuance of medical certificates can be time-consuming for both staff and students. Automating this process through the PIMS can streamline operations and improve student satisfaction \cite{hu2023}.

\subsection{Information Dissemination}
University clinics play a crucial role in health promotion and disease prevention. However, disseminating health information effectively across a large and dispersed campus can be challenging. Leveraging the PIMS to deliver targeted health campaigns, educational materials, and reminders can enhance the clinic’s reach and impact \cite{iribarren2021}.

\subsection{Patient Feedback and Evaluation}
Collecting and analyzing patient feedback is essential for continuous improvement in healthcare settings. Electronic feedback systems integrated into the PIMS can streamline data collection, increase response rates, and enable the clinic to make data-driven decisions to improve service quality \cite{ferreira2023}.

\section{Latest Trends in HIS Technology}

Modernizing the USC Clinic’s PIS requires consideration of the latest technological advancements:

\begin{enumerate}
    \item \textbf{Cloud-based Solutions:} Cloud-based HIS platforms offer scalability, cost-effectiveness, and remote accessibility, making them an attractive option for university clinics \cite{ali2023}.
    \item \textbf{Mobile Health (mHealth):} mHealth applications can empower patients to manage their health information, schedule appointments, and communicate with providers, enhancing patient engagement and convenience \cite{free2013}.
    \item \textbf{Artificial Intelligence (AI):} AI and machine learning can be integrated into HIS to improve diagnostics, predict patient outcomes, and personalize treatment plans, ultimately enhancing the quality of care \cite{ching2018}.
\end{enumerate}
